\documentclass[oneside,a4paper,14pt]{extarticle} %размер шрифта 14
\usepackage[T1,T2A]{fontenc}
\usepackage[
	a4paper,
	letterpaper,
	top=2cm,
	bottom=2cm,
	left=2.5cm,
	right=1.5cm
]{geometry}
\usepackage[utf8]{inputenc} %кодировка текста
\usepackage[russian]{babel} %поддержка русского языка
\usepackage{textcomp} %текстовые символы
\usepackage{indentfirst} %корректировка отступов
\usepackage{listings} %для кода !!
\usepackage{xcolor}
\usepackage{enumitem} %для нумерации буквами
\usepackage{graphicx} %работа с изображениям
\usepackage{mwe} % for blindtext and example-image-a in example
\usepackage{wrapfig}
\usepackage{caption}
\usepackage{amsmath}  % для формул и символов
\usepackage{amsfonts}
\usepackage{float}
\usepackage{amsthm}
\usepackage{graphicx}
\usepackage[all]{xy}
\usepackage[breaklinks]{hyperref}
%%размеркегляузаголоковразделов
\usepackage{titlesec}
\titleformat{\section}
{\normalsize\bfseries}
{\thesection} {1em}{}
\titleformat{\subsection}
{\normalsize\bfseries}
{\thesubsection} {1em}{}
\titleformat{\subsubsection}
{\normalsize\bfseries}
{\thesubsection} {1em}{}

\renewcommand\baselinestretch{1.33}\normalsize % межстрочный интервал
\setlength{\parindent}{1.25cm}
\usepackage{indentfirst}

% настройки листинга
\lstset{
  language=Python,
  basicstyle=\ttfamily\small,
  keywordstyle=\color{blue},
  commentstyle=\color{green!60!black},
  stringstyle=\color{red},
  numbers=left,
  numberstyle=\tiny\color{gray},
  frame=single,
  breaklines=true,
  showstringspaces=false,
  tabsize=2,
  inputencoding=cp1251,
  extendedchars=true
}

\begin{document}
	\newpage\thispagestyle{empty}
	\begin{center}
		МИНИСТЕРСТВО НАУКИ И ВЫСШЕГО ОБРАЗОВАНИЯ\\
			РОССИЙСКОЙ ФЕДЕРАЦИИ
			ФЕДЕРАЛЬНОЕ ГОСУДАРСТВЕННОЕ БЮДЖЕТНОЕ\\
			ОБРАЗОВАТЕЛЬНОЕ
			УЧРЕЖДЕНИЕ ВЫСШЕГО ОБРАЗОВАНИЯ\\
			«ВЯТСКИЙ ГОСУДАРСТВЕННЫЙ УНИВЕРСИТЕТ»\\
			Институт математики и информационных систем\\
			Факультет автоматики и вычислительной техники\\
			Кафедра электронных вычислительных машин
	\end{center}
	\vspace{20mm}

	\begin{center}
		Отчёт по лабораторной работе №7\\
		по дисциплине\\
		<<Программирование>>\\
	\end{center}
	\vspace{48mm}

	\noindent
	Выполнил студент гр. ИВТб-1303-06-00 \hfill \rule[-0,5mm]{30mm}{0.15mm}\,/Гортоломей И.К./
	\newline
	Проверил преподаватель кафедры ЭВМ \hfill \rule[-0,5mm]{30mm}{0.15mm}\,/Баташев П.А./

	\vfill
	\begin{center}
		Киров\\
		2026
	\end{center}

	\newpage
	\thispagestyle{empty}

% ==================== 1. ЦЕЛЬ И ЗАДАЧИ ====================
\section{Цель и задачи работы}
\textbf{Цель работы:} освоение практических навыков создания интерактивных дашбордов в Python, включая интеграцию \texttt{pandas DataFrame} с графическими библиотеками \texttt{matplotlib} и \texttt{seaborn}, встраивание графиков в окно \texttt{Tkinter} через бэкенд \texttt{FigureCanvasTkAgg}, реализацию событийно-ориентированной архитектуры (callbacks, event loop), динамическую фильтрацию, агрегацию и feature engineering <<на лету>>, а также проектирование устойчивого процедурного кода без использования классов.

\textbf{Задачи:}
\begin{enumerate}
    \item Сформировать и подготовить набор данных с категорийными и временными признаками.
    \item Реализовать предобработку данных: очистку выбросов, клаппинг физических величин, расчёт скользящих средних и производных признаков.
    \item Интегрировать графический движок Matplotlib/Seaborn с GUI Tkinter через адаптер \texttt{FigureCanvasTkAgg}.
    \item Разработать панель управления с интерактивными элементами (выпадающие списки, переключатели агрегации, кнопки смены типа графика).
    \item Реализовать механизм безопасного обновления данных через \texttt{.copy()} и перерисовки через \texttt{canvas.draw\_idle()} без блокировки основного потока.
    \item Внедрить экспорт визуализаций в PNG/PDF.
\end{enumerate}

\newpage
% ==================== 2. ОПИСАНИЕ ВАРИАНТА ====================
\section{Описание варианта}
\textbf{Вариант №20: Водоподготовка}\\
\textbf{Предметная область:} Мониторинг датчиков очистных сооружений. Контроль гарантирует соответствие санитарным нормам качества питьевой воды.\\
\textbf{Основные метрики (поля датасета):}
\begin{itemize}
    \item \texttt{ts} (int32) -- время замера (Unix timestamp);
    \item \texttt{filter\_id} (int16) -- идентификатор фильтра;
    \item \texttt{turb} (float32) -- турбидность, NTU;
    \item \texttt{ph} (float32) -- кислотность воды;
    \item \texttt{cl} (float32) -- остаточный хлор, мг/л;
    \item \texttt{flow} (float32) -- расход воды, м\textsuperscript{3}/ч.
\end{itemize}

\textbf{Условия фильтрации и очистки:}
\begin{itemize}
    \item \texttt{turb < 0} $\rightarrow$ 0; \texttt{cl < 0} $\rightarrow$ 0 (физическая невозможность отрицательных значений);
    \item \texttt{ph} обрезается (clip) по диапазону $[6.5; 8.5]$ согласно санитарным нормам;
    \item Выбросы по турбидности заменяются медианой группы (IQR-метод по \texttt{filter\_id});
    \item Редкие фильтры ($<1\%$ от общего объёма) агрегируются в категорию \texttt{0}.
\end{itemize}

\textbf{Группировка и оконные функции:}
\begin{itemize}
    \item Группировка по \texttt{filter\_id}: расчёт среднего \texttt{turb} (эффективность очистки) и максимума \texttt{flow} (нагрузка на систему);
    \item Скользящее окно $k=35$: расчёт \texttt{turb\_ma} для выявления долгосрочных трендов;
    \item \texttt{np.diff(flow)} $\rightarrow$ изменение расхода между замерами (индикатор работы насосов или засорения).
\end{itemize}

\textbf{Акцент аналитики:} Обратная зависимость чистоты воды и турбидности. При фильтрации аномалий каждый параметр обрабатывается независимо, без взаимного влияния масок.\\
\textbf{Типы графиков:} Линейный (динамика по неделям), столбчатый (сезонная агрегация), точечный (pH vs turbidity с цветовым кодированием статуса), гистограмма (распределение турбидности).

% \newpage
% ==================== 3. ЛИСТИНГ КОДА ====================
\section{Листинг программного кода}
Код реализован в процедурном стиле без использования классов. Логика разбита на этапы в соответствии с методическими указаниями.

\lstinputlisting[caption={Этап 0-1: Инициализация и загрузка данных}]{1.nonutf.txt}
\newpage

\lstinputlisting[caption={Этап 2: Предобработка и Feature Engineering}]{2.nonutf.txt}


\begin{lstlisting}[caption={Этап 3: Встраивание Figure в Tkinter}]
plot_frame = tk.Frame(root, bg="white", relief=tk.SUNKEN, bd=1)
plot_frame.pack(fill=tk.BOTH, expand=True, padx=10, pady=5)

fig = plt.Figure(figsize=(10, 6), dpi=100)
canvas = FigureCanvasTkAgg(fig, master=plot_frame)
canvas.get_tk_widget().pack(fill=tk.BOTH, expand=True)

toolbar = NavigationToolbar2Tk(canvas, plot_frame)
toolbar.update()
toolbar.pack(side=tk.TOP, fill=tk.X)
\end{lstlisting}
\newpage

\lstinputlisting[caption={Этап 4: Функции отрисовки графиков (Seaborn/Matplotlib)}]{4.nonutf.txt}


\lstinputlisting[caption={Этап 5-6: Панель управления и интерактивность}]{5.nonutf.txt}


% \newpage
% ==================== 4. СКРИНШОТЫ РАБОТЫ ====================
\section{Скриншоты работы программы}
\begin{figure}[H]
    \centering
    \includegraphics[width=16cm]{1.png}
    \caption{Линейный график: динамика турбидности по всем фильтрам}
\end{figure}
\newpage
\begin{figure}[H]
    \centering
    \includegraphics[width=17cm]{2.png}
    \caption{Линейный график: динамика турбидности по одному фильтру}
\end{figure}
% \newpage
\begin{figure}[H]
    \centering
    \includegraphics[width=17cm]{3.png}
    \caption{Столбчатая диаграмма: агрегация турбидности по сезонам}
\end{figure}
\newpage
\begin{figure}[H]
    \centering
    \includegraphics[width=17cm]{4.png}
    \caption{Точечная диаграмма: зависимость turb от pH с цветовым кодированием статуса}
\end{figure}
% \newpage
\begin{figure}[H]
    \centering
    \includegraphics[width=13cm]{5.png}
    \caption{Диалог сохранения в PNG/PDF}
\end{figure}

\newpage
% ==================== 5. АНАЛИЗ РЕЗУЛЬТАТОВ ====================
\section{Анализ результатов}
В ходе выполнения лабораторной работы были получены следующие аналитические инсайты и выявлены системные ограничения:

\textbf{Полученные инсайты:}
\begin{itemize}
    \item Применение скользящего среднего (\texttt{window=35}) эффективно сглаживает шумовые выбросы турбидности, позволяя чётко выделить долгосрочные тренды эффективности фильтрации. Резкие пики на сырых данных соответствуют кратковременным перегрузкам насосов, а не системному ухудшению качества воды.
    \item Независимая обработка масок аномалий (IQR для \texttt{turb}, clip для \texttt{ph}) предотвращает каскадное искажение данных. Цветовое кодирование \texttt{ph\_status} на scatter-plot наглядно демонстрирует отсутствие прямой линейной корреляции между кислотностью и мутностью, что подтверждает необходимость раздельного контроля параметров.
    \item Агрегация по сезонам выявляет, что максимальные значения турбидности приходятся на весенний период (паводковые нагрузки), что коррелирует с увеличением \texttt{flow}.
\end{itemize}

\textbf{Выявленные ограничения:}
\begin{itemize}
    \item Архитектура Tkinter является однопоточной. При обработке датасетов свыше $10^7$ строк пересчёт \texttt{groupby()} и отрисовка Seaborn могут вызывать кратковременные <<фризы>> интерфейса.
    \item Статический рендеринг Matplotlib/Seaborn не поддерживает аппаратное ускорение GPU, что ограничивает частоту обновления в режимах реального времени (Real-time streaming).
    \item Использование \texttt{canvas.draw\_idle()} ставит задачу в очередь событий, но не гарантирует мгновенного отражения данных при частых изменениях ползунков или фильтров.
\end{itemize}

\textbf{Пути масштабирования решения:}
\begin{enumerate}
    \item Перенос тяжёлых вычислений в фоновый поток с использованием \texttt{root.after()} или \texttt{queue.Queue} для безопасной передачи результатов в GUI-поток.
    \item Замена Seaborn на \texttt{Plotly} или \texttt{Bokeh} для нативной поддержки WebGL-ускорения и интерактивного масштабирования без потери производительности.
    \item Интеграция \texttt{Dask DataFrame} или \texttt{Polars} для распределённой обработки данных, выходящих за пределы RAM, с сохранением логики фильтрации на уровне API.
\end{enumerate}

% \newpage
% ==================== 6. ВЫВОДЫ ====================
\section{Выводы}
В ходе выполнения лабораторной работы были успешно освоены ключевые технологии создания интерактивных аналитических панелей в экосистеме Python:
\begin{itemize}
    \item Реализована событийно-ориентированная архитектура с использованием callback-функций и цикла событий \texttt{mainloop()}, что обеспечило мгновенный отклик интерфейса на действия пользователя.
    \item Освоено встраивание графических объектов Matplotlib в оконную систему Tkinter через адаптер \texttt{FigureCanvasTkAgg}, включая корректную работу панели навигации (\texttt{NavigationToolbar2Tk}).
    \item На практике применены методы динамической фильтрации (\texttt{df.loc[]}), агрегации (\texttt{groupby().agg()}), оконных функций (\texttt{rolling()}) и биннинга (\texttt{pd.cut()}) библиотеки pandas. Соблюдено правило обязательного копирования данных (\texttt{.copy()}) для предотвращения \texttt{SettingWithCopyWarning} и искажения исходного датасета.
    \item Разработано устойчивое процедурное приложение без использования ООП-паттернов, что упростило отладку и чтение кода.
\end{itemize}

\textbf{Трудности:} Основным вызовом стала синхронизация состояния данных между этапом предобработки и этапом отрисовки. Некорректное использование ссылок на DataFrame приводило к накоплению артефактов на графиках. Проблема была решена строгим разделением \texttt{df\_processed} (статический эталон) и \texttt{df\_work} (рабочая копия), а также использованием \texttt{fig.clear()} перед каждой отрисовкой.\\
\textbf{Пути улучшения:} Добавление логгирования событий, внедрение прогресс-баров для длительных вычислений, поддержка темной темы интерфейса и автоматическое сохранение состояния фильтров между сессиями.

% \newpage
% ==================== 7. ССЫЛКА НА РЕПОЗИТОРИЙ ====================
\section{Ссылка на репозиторий}
Исходный код проекта, скрипты генерации данных и инструкции по запуску доступны по ссылке:\\
\url{https://github.com/GIKExe/Subject/tree/main/Proga/laba7}

\end{document}